\documentclass[11pt,a4paper]{article}
\usepackage[utf8]{inputenc}
\usepackage{amsmath,amssymb,amsthm}
\usepackage{mathtools}
\usepackage{geometry}
\usepackage{hyperref}
\usepackage{enumitem}
\usepackage{xcolor}

\geometry{margin=2.5cm}

\newtheorem{theorem}{Theorem}[section]
\newtheorem{lemma}[theorem]{Lemma}
\newtheorem{proposition}[theorem]{Proposition}
\newtheorem{definition}[theorem]{Definition}
\newtheorem{remark}[theorem]{Remark}

\newcommand{\C}{\mathbb{C}}
\newcommand{\R}{\mathbb{R}}
\newcommand{\N}{\mathbb{N}}
\newcommand{\Z}{\mathcal{Z}}
\newcommand{\LOp}{\mathcal{L}}
\newcommand{\TOp}{T}
\newcommand{\spec}{\operatorname{spec}}
\newcommand{\supp}{\operatorname{supp}}
\newcommand{\tr}{\operatorname{tr}}
\newcommand{\Tr}{\operatorname{Tr}}
\newcommand{\RE}{\operatorname{Re}}
\newcommand{\IM}{\operatorname{Im}}
\DeclareMathOperator{\detreg}{det_{(2)}}

\title{\bf Research Programme:\\
Analytic Estimates for the Greedy Harmonic / Iota Walk\\
Approach to the Riemann Hypothesis}
\author{}
\date{\today}

\begin{document}
\maketitle

\begin{abstract}
This document outlines a systematic programme of research aimed at providing the missing analytic estimates required to turn the geometric formulations of the Riemann Hypothesis—via the greedy harmonic expansion, the iota walk, and the associated Jacobi matrix—into a rigorous proof.  The programme does not claim to have such a proof; it sets out the concrete mathematical problems that must be solved, the techniques that may be employed, and a realistic assessment of the obstacles.
\end{abstract}

\section{Introduction}
The Riemann Hypothesis (RH) states that every non‑trivial zero of the zeta function lies on the critical line $\operatorname{Re}(s)=\frac12$.  Recent work has produced several equivalent formulations of the hypothesis in a geometric, explicitly real‑analytic language:

\begin{enumerate}[nosep]
    \item The \textbf{greedy harmonic expansion} expresses $\sin\theta$ as a limit of sums involving greedy binary digits $\delta_n(\theta)$.
    \item The \textbf{iota walk} views the Dirichlet eta series $\eta(s)=\sum_{n=1}^\infty (-1)^{n-1}n^{-s}$ as a planar walk, whose steps are $v_n(s)=(-1)^{n-1}n^{-s}$, and introduces a \textbf{greedy condition} (GC) $ \langle\mathbf{S}_{N-1},v_N\rangle \le 0$ that is equivalent to $\sigma=1/2$ for zeros.
    \item A \textbf{continuous analogue} of the walk is derived from Riemann's integral representation, leading to a monotonicity condition on the distance $|S(\alpha)|$.
    \item A \textbf{Jacobi matrix} $J(t)$ can be constructed from the squared distances of the iota partial sums, with the property that $0\in\operatorname{spec}(J(\gamma))$ exactly when $\eta(\frac12+i\gamma)=0$.
\end{enumerate}

These formulations are mathematically rigorous, but they merely restate the Riemann Hypothesis in a different language.  The crucial gap that remains is the lack of \textbf{analytic estimates} that would allow one to prove, from the structural properties of the greedy digits or the iota walk, that the only possible zeros are those on the critical line.

This research programme identifies the precise estimates that are missing and proposes a route to obtaining them.  It is \textbf{not} a proof; it is a roadmap for the necessary mathematical work.

\section{The Core Objects and the Missing Link}
\label{sec:core}

We recall the central objects:

\begin{itemize}
    \item The greedy digits $\delta_n(\theta)$ (binary, piecewise constant on intervals whose number grows super‑exponentially).
    \item The iota walk $S_N(t)=\sum_{n=1}^N (-1)^{n-1}n^{-1/2}e^{-it\ln n}$ (real and imaginary parts $X_N,Y_N$).
    \item The greedy condition: for a zero $\gamma$, for all $N$, $X_{N-1}x_N+Y_{N-1}y_N\le 0$.
    \item The continuous analogue $S(\alpha) = 2\int_0^\alpha \frac{\sin((\frac12+it)\theta)}{(\cos\theta)^2(e^{2\pi\tan\theta}-1)}\,d\theta -\frac32$, with the monotone condition $\operatorname{Re}(\overline{S(\alpha)}S'(\alpha))\le0$.
    \item The Jacobi matrix $J(t)$ built from moments $m_k(t)=\sum_{N=1}^\infty N^k |S_N(t)|^2 \rho_N$ (with appropriate regularisation).
\end{itemize}

The missing link is an \textbf{analytic inequality} that, when applied to the iota walk or the continuous walk, forces $\sigma$ to be $1/2$ whenever the walk converges to zero \emph{and} satisfies the greedy condition (or its continuous analogue).  Equivalently, one must prove that for $\sigma\neq 1/2$, either the walk cannot converge to zero, or the greedy condition must fail at some finite step.

\section{Objective 1: Asymptotic Estimates for the Iota Walk}
\label{sec:asymp-iota}

The first block of estimates concerns the behaviour of the partial sums $S_N(t)$ for large $N$ when $\sigma\neq 1/2$.  Let $s=\sigma+it$.

\subsection{Sub‑problem 1a: Upper and lower bounds for $|S_N|$}
Prove that for $\sigma>1/2$, the partial sums converge absolutely and the limit is non‑zero (except perhaps for finitely many $t$), and obtain an effective lower bound $|S_N| \ge c N^{1/2-\sigma}$ for sufficiently large $N$.  For $\sigma<1/2$, show that $|S_N|$ grows at least like $N^{1/2-\sigma}$, and hence the walk cannot tend to zero.

\subsection{Sub‑problem 1b: Oscillation control and failure of GC}
Show that for $\sigma<1/2$, the greedy condition must be violated at infinitely many $N$.  This would follow if one can prove that the step vectors $v_N$ eventually point in the direction of the current sum, i.e.\ $\langle S_{N-1},v_N\rangle >0$ eventually.  This requires a detailed estimate of the correlation between $S_{N-1}$ and $v_N$, which depends on the distribution of the angles $t\ln n$ modulo $2\pi$.

\subsection{Methodology}
\begin{itemize}
    \item Use the classical theory of exponential sums (van der Corput, Weyl) to estimate the sums $\sum_{n\le N} n^{-s}$.
    \item Combine with partial summation to handle the alternating signs.
    \item For GC, one may attempt to link $\langle S_{N-1},v_N\rangle$ to the derivative of $|S_{N-1}|^2$, since $|S_N|^2-|S_{N-1}|^2 = 2\langle S_{N-1},v_N\rangle+|v_N|^2$.
    \item Known results on the distribution of $\ln n$ modulo 1 (the “logarithm circle problem”) may be crucial.
\end{itemize}

\section{Objective 2: Estimates for the Continuous Walk}
\label{sec:continuous}

The continuous greedy condition is $\frac{d}{d\alpha}|S(\alpha)|^2 \le 0$.  This is a differential inequality.  The objective is to prove that this inequality \emph{uniquely} determines the critical line.

\subsection{Sub‑problem 2a: Expression via the harmonic sine expansion}
Using the harmonic sine expansion,
\[
S(\alpha) = -\frac32 + 2\sum_{m=0}^\infty a_m \int_0^\alpha \frac{\theta^{2m+1}\cosh(t\theta)}{(\cos\theta)^2(e^{2\pi\tan\theta}-1)}\,d\theta
+ 2i\sum_{m=0}^\infty b_m \int_0^\alpha \frac{\theta^{2m}\sinh(t\theta)}{(\cos\theta)^2(e^{2\pi\tan\theta}-1)}\,d\theta .
\]
One can rewrite the derivative $|S(\alpha)|^2' = 2\operatorname{Re}(\overline{S(\alpha)}S'(\alpha))$ as a double series.  The task is to show that this double series is non‑positive for all $\alpha\in[0,\pi/2]$ \emph{only when} $\sigma=1/2$ (note: the representation above is for $\sigma=1/2$ already; for general $\sigma$ one must use $\sin((\sigma+it)\theta)$ which introduces $\sin(\sigma\theta)$ and $\cos(\sigma\theta)$).  Thus we must extend the harmonic sine expansion to arbitrary $\sigma$ using the greedy digits $\delta_n(\theta)$ to represent $\sin(\sigma\theta)$.

\subsection{Sub‑problem 2b: Monotonicity impossibility off the line}
Prove that for $\sigma>1/2$, the distance $|S(\alpha)|$ must eventually increase because the hyperbolic decay of the integrand is too weak to bring the curve back to zero while keeping the direction condition.  For $\sigma<1/2$, the integrand grows and the distance must increase.  This would show that the continuous greedy condition can hold only when $\sigma=1/2$.

\subsection{Methodology}
\begin{itemize}
    \item Variation of parameters: $|S(\alpha)|^2$ satisfies a differential equation; analyse its sign near $\alpha=0$ and near the end point.
    \item Use integral representations of the functions involved to obtain explicit formulas for $|S(\alpha)|^2'$.
    \item Exploit the super‑exponential growth of the selected indices to deduce that the series for $\sin(\sigma\theta)$ heavily weights certain indices, creating an imbalance that violates monotonicity unless $\sigma=1/2$.
\end{itemize}

\section{Objective 3: Spectral Estimates for the Jacobi Matrix}
\label{sec:jacobi-est}

The Jacobi matrix $J(t)$ is defined from the moments $m_k(t)=\sum_{N=1}^\infty N^k |S_N(t)|^2 e^{-\epsilon\ln^2 N}$.  The zero eigenvalue appears iff $S_N(t)\to0$.

\subsection{Sub‑problem 3a: Regularisation independence}
Show that the asymptotic presence of a zero eigenvalue (or spectral gap at zero) is independent of the regularisation choice, provided it is sufficiently mild.

\subsection{Sub‑problem 3b: Connection to zeros on the line}
Prove that if $0\in\operatorname{spec}(J(t))$, then necessarily $t$ corresponds to a zero of $\eta(\frac12+it)$.  This is fairly straightforward from the moment inversion, but one must control the tails.

\subsection{Sub‑problem 3c: Restriction to the critical line}
Show that if $t$ is such that $0\in\operatorname{spec}(J(t))$, then the associated $s=\frac12+it$ must satisfy $\operatorname{Re}(s)=1/2$.  This is the heart of the matter.  One would need to prove that for $\sigma\neq 1/2$, the moments $m_k(t)$ grow or decay at rates incompatible with having a zero eigenvalue.  This reduces to showing that for $\sigma\neq 1/2$, the partial sums $S_N(t)$ do not tend to zero, which is Objective 1.

Thus Objective 3 is essentially equivalent to proving that the iota walk cannot converge to zero off the critical line.  The Jacobi matrix just repackages the same problem.

\section{Objective 4: The Greedy Digits and the GDST Operator}
\label{sec:greedy-op}

Although the GDST programme failed due to a false identification, the correlation operator $T_\sigma$ and the Fredholm determinant identity remain valid.  

\subsection{Sub‑problem 4a: Estimate the eigenvalues of $T_\sigma$}
Prove that the eigenvalues $\lambda_i(\sigma)$ decay like $i^{-2}$ and that the trace identities hold.  This is already largely established.

\subsection{Sub‑problem 4b: Relate $T_\sigma$ to the iota walk}
Find a direct relationship between the correlation kernel $K_\sigma(\theta,\phi)$ and the iota walk, possibly through the formula $E(\theta,s)=\sum_{n:\delta_n(\theta)=1}(n+2)^{-s}$ which appears in the splitting of $\zeta(s)$.  The hope is that the iota walk’s greedy condition translates into a positivity property of the operator $(1-T_\sigma)^{-1}$, allowing an alternative spectral proof.

\subsection{Methodology}
\begin{itemize}
    \item Express the partial sums $S_N(t)$ as a linear combination of the eigenvectors of $T_{1/2}$.
    \item The GC may then be interpreted as a condition on the coefficients of the expansion, leading to an inequality that forces the spectral measure to be supported on $[0,1]$, which in turn forces the zero to lie on the line.
\end{itemize}

\section{Planned Activities and Timeline}
\label{sec:timeline}

\begin{enumerate}[nosep]
    \item \textbf{Months 1‑6:} Develop the asymptotic estimates for the iota walk (Objective 1).  Focus on the case $\sigma>1/2$ first, using classical exponential sum techniques.  Aim for a proof that $S_N(t)\to L\neq0$ for $\sigma>1/2$ (except possibly at $t$ corresponding to zeros on $\sigma=1/2$).
    \item \textbf{Months 7‑12:} Extend the estimates to $\sigma<1/2$, proving divergence or violation of GC.  Concurrently, work on the continuous walk (Objective 2) to prove the monotonicity condition characterises $\sigma=1/2$.
    \item \textbf{Months 13‑18:} Investigate the Jacobi matrix (Objective 3) and its regularisation.  Attempt to prove that zero eigenvalue implies $\sigma=1/2$ via the moment estimates.
    \item \textbf{Months 19‑24:} Explore the connection between the GDST operator and the iota walk (Objective 4).  Attempt to re‑formulate the failed GDST argument using the iota greedy condition to avoid the earlier contradiction.
    \item \textbf{Ongoing:} Computational experiments to test the emerging inequalities for high zeros and to guide the analytic estimates.
\end{enumerate}

\section{Assessment of Feasibility}
\label{sec:feasibility}

The programme is ambitious and likely to encounter severe technical difficulties:

\begin{itemize}
    \item The exponential sums involved are notoriously delicate; controlling the oscillations for all $N$ and $t$ simultaneously is a major challenge.
    \item The greedy condition involves dot products of partial sums, which are sums of products of oscillatory terms; obtaining a uniform sign for all $N$ seems extremely difficult.
    \item The continuous greedy condition is a global differential inequality; proving it characterises a specific parameter is a problem of a kind that rarely has a clean solution.
    \item The Jacobi matrix approach ultimately reduces to the same convergence/divergence estimates as Objective 1.
\end{itemize}

Nevertheless, even a partial success—e.g., proving that for $\sigma>1/2$ the walk cannot converge to zero—would eliminate half of the critical strip and represent significant progress towards RH.  Moreover, the geometric reformulations may inspire new types of estimates (e.g., using the theory of orthogonal polynomials, martingales, or potential theory) that have not yet been explored.

\section{Conclusion}
This research programme outlines a concrete path to supply the missing analytic estimates that currently prevent the greedy harmonic/iota walk approach from becoming a proof of the Riemann Hypothesis.  The estimates are of a classical nature—exponential sums, differential inequalities, moment problems—but they are pushed into a regime where current technology is insufficient.  The programme is therefore both a call for new ideas and a systematic attempt to apply existing methods to a well‑defined set of problems.

\begin{thebibliography}{9}
\bibitem{GDST} V.\ Geere et al., \textit{The Greedy Dirichlet Sub‑Sum Transform and the Riemann zeta function}, preprint, 2026.
\bibitem{Iota} V.\ Geere, \textit{On the Iota Function and the Greedy Condition: A Reformulation, Not a Proof}, manuscript, 2026.
\bibitem{Titchmarsh} E.\,C.\ Titchmarsh, \textit{The Theory of the Riemann Zeta‑function}, 2nd ed., Oxford Univ.\ Press, 1986.
\bibitem{Mayer} D.\ Mayer, \textit{Continued fractions and related transformations}, in: \textit{Ergodic Theory, Symbolic Dynamics, and Hyperbolic Spaces}, Oxford Univ.\ Press, 1991.
\end{thebibliography}

\end{document}